\documentclass[13pt,xcolor={usenames,dvipsnames}]{beamer}
%% \mode<presentation>
%\usetheme{Warsaw}
\usetheme[height=8mm]{Rochester}
\setbeamertemplate{blocks}[rounded][shadow=true]
\setbeamertemplate{navigation symbols}{}

% if add some extra slides  -- Jiguo
\newif\ifnoextra
\noextratrue       % this will exclude all extra sides
% \noextrafalse      % this will include all extra sides

% if show page number    -- Jiguo
\newif\ifpgnum
\pgnumtrue       % this will exclude all extra sides
% \pgnumfalse      % this will include all extra sides

\ifpgnum
\setbeamertemplate{footline}[page number]
\else 
\fi


\usepackage[absolute,overlay]{textpos}
\setlength{\TPHorizModule}{1mm}
\setlength{\TPVertModule}{1mm}
\usepackage[absolute,overlay]{textpos}
\usepackage[utf8]{inputenc}
\usepackage{hyperref}
\usepackage{color}
\usepackage{soul}
\usepackage[T1]{fontenc}
\usepackage[utf8]{inputenc}
\usepackage{environ}

\usepackage{tikz}
\newcommand*\circled[1]{\tikz[baseline=(char.base)]{
            \node[shape=circle,draw,inner sep=2pt] (char) {#1};}}
%\usepackage{charter}

%\beamertemplateballitem
%\setbeameroption{show notes}
%\usecolortheme{lily}
%\usecolortheme{whale}


%\usepackage[usenames,dvipsnames]{xcolor}
\usepackage{tikz}
\usetikzlibrary{calc,matrix}

\definecolor{light-gray}{gray}{0.8}
\definecolor{dgray}{rgb}{0.3,0.3,0.5}
\usecolortheme[named=dgray]{structure}
\definecolor{grey}{rgb}{0.3,0.3,0.5}
\definecolor{darkgreen}{rgb}{0.0, 0.2, 0.13}
\definecolor{bole}{rgb}{0.47, 0.27, 0.23}
\definecolor{bittersweet}{rgb}{1.0, 0.44, 0.37}
\definecolor{indiagreen}{rgb}{0.07, 0.53, 0.03}

% code by Andrew:
% http://tex.stackexchange.com/a/28452/13304
\makeatletter
\let\matamp=&
\catcode`\&=13
\makeatletter
\def&{\iftikz@is@matrix
  \pgfmatrixnextcell
  \else
  \matamp
  \fi}
\makeatother

\newcounter{lines}
\def\endlr{\stepcounter{lines}\\}

\newcounter{vtml}
\setcounter{vtml}{0}

\newif\ifvtimelinetitle
\newif\ifvtimebottomline
\tikzset{description/.style={
    column 2/.append style={#1}
  },
  timeline color/.store in=\vtmlcolor,
  timeline color=red!255!black,
  timeline color st/.style={fill=\vtmlcolor,draw=\vtmlcolor},
  use timeline header/.is if=vtimelinetitle,
  use timeline header=false,
  add bottom line/.is if=vtimebottomline,
  add bottom line=false,
  timeline title/.store in=\vtimelinetitle,
  timeline title={},
  line offset/.store in=\lineoffset,
  line offset=4pt,
}

\NewEnviron{vtimeline}[1][]{%
  \setcounter{lines}{1}%
  \stepcounter{vtml}%
  \begin{tikzpicture}[column 1/.style={anchor=east},
    column 2/.style={anchor=west},
    text depth=0pt,text height=1ex,
    row sep=1ex,
    column sep=1em,
    #1
    ]
    \matrix(vtimeline\thevtml)[matrix of nodes]{\BODY};
    \pgfmathtruncatemacro\endmtx{\thelines-1}
    \path[timeline color st] 
    ($(vtimeline\thevtml-1-1.north east)!0.5!(vtimeline\thevtml-1-2.north
    west)$)--
    ($(vtimeline\thevtml-\endmtx-1.south
    east)!0.5!(vtimeline\thevtml-\endmtx-2.south west)$);
    \foreach \x in {1,...,\endmtx}{
      \node[circle,timeline color st, inner sep=0.25pt, draw=white, thick] 
      (vtimeline\thevtml-c-\x) at 
      ($(vtimeline\thevtml-\x-1.east)!0.5!(vtimeline\thevtml-\x-2.west)$){};
      \draw[timeline color st](vtimeline\thevtml-c-\x.west)--++(-3pt,0);
    }
    \ifvtimelinetitle%
    \draw[timeline color st]([yshift=\lineoffset]vtimeline\thevtml.north
    west)--
    ([yshift=\lineoffset]vtimeline\thevtml.north east);
    \node[anchor=west,yshift=16pt,font=\large]
    at (vtimeline\thevtml-1-1.north west) 
    {\textsc{Timeline} \hspace{1cm} \textit{\vtimelinetitle}};
    \else%
    \relax%
    \fi%
    \ifvtimebottomline%
    \draw[timeline color
    st]([yshift=-\lineoffset]vtimeline\thevtml.south west)--
    ([yshift=-\lineoffset]vtimeline\thevtml.south east);
    \else%
    \relax%
    \fi%
  \end{tikzpicture}
}


\usepackage{listings}
\lstset{numbers=none,language=[ISO]C++,tabsize=4,
  frame=single,
  basicstyle=\small,
  showspaces=false,showstringspaces=false,
  showtabs=false,
  keywordstyle=\color{blue}\bfseries,
  commentstyle=\color{red},
}

\usepackage[absolute,overlay]{textpos}
\usepackage{tikz}

% We need lots of libraries...
\usetikzlibrary{%
  arrows,%
  calc,%
  fit,%
  patterns,%
  plotmarks,%
  shapes.geometric,%
  shapes.misc,%
  shapes.symbols,%
  shapes.arrows,%
  shapes.callouts,%
  shapes.multipart,%
  shapes.gates.logic.US,%
  shapes.gates.logic.IEC,%
  er,%
  automata,%
  backgrounds,%
  chains,%
  topaths,%
  trees,%
  petri,%
  mindmap,%
  matrix,%
  calendar,%
  folding,%
  fadings,%
  through,%
  positioning,%
  scopes,%
  decorations.fractals,%
  decorations.shapes,%
  decorations.text,%
  decorations.pathmorphing,%
  decorations.pathreplacing,%
  decorations.footprints,%
  decorations.markings,%
  shadows} 

\newcommand{\specialcell}[2][c]{%
  \begin{tabular}[#1]{@{}c@{}}#2\end{tabular}}

%**
% \PutAt<overlay spec>[<box width>]{(<x>, <y>)}{<content>}
%
% real absolute positioning of <content> on a slide, if content is a
% figure,
% minipage or whatever kind of LR-box, the <box width> argument may be
% omitted
%
%
% implementation notes: 
%   - based on   \usepackage[absolute,overlay]{textpos}
%   - NOT combinable with any beamer feature that is based on pgfpages
%     (such as dual-screen support, built-in 2up handouts, etc.), as
% textpos 
%     and pgfpates interfere at the shippout-level.
%


%def\Put(#1,#2)#3{\leavevmode\makebox(0,0){\put(#1,#2){#3}}}
\newcommand<>{\PutJ}[3][0pt]{%
  {\leavevmode\makebox(0,0)
    {
      \put(#1,#2){#3}
    }
  }
}

\newcommand<>{\PutAt}[3][0pt]{%
  {\only#4{\begin{textblock*}{#1}#2%
        #3
      \end{textblock*}}}%
}

% **
% \ShowPutAtGrid
% 
% draws a helpful grid on the current slide to figure <x> and <y>
% parameters for \PutAt
% 
\newcommand{\ShowPutAtGrid}{
  \begin{textblock*}{128mm}(0cm,0cm)
    \tikz[color=red!20!white]\draw[very thin, step=5mm] (0mm,0mm) grid
    (130mm,100mm);
  \end{textblock*}
  \begin{textblock*}{128mm}(0cm,0cm)
    \begin{tikzpicture}[color=red]
      \draw[step=1cm] (0,0mm) grid (130mm,100mm);   
      \foreach \n in {0,...,12}
      \draw[xshift=.5mm,yshift=-1.5mm, inner sep=0pt, anchor=west]
      (\n,10) node {\scriptsize{\textbf{\n}}};
      \foreach \n in {1,...,9}
      \draw[xshift=.5mm,yshift=-1.5mm, inner sep=0pt, anchor=west]
      (0,10-\n) node {\scriptsize{\textbf{\n}}};
    \end{tikzpicture}
  \end{textblock*}
}

% an box that does nothing
\newcommand<>{\WhiteBox}[2][]{%
  \only#3{\tikz[#1, every
    node/.style={shape=rectangle,fill=white, #1}]\node []{#2};}
}

% **
% \NormalBox<overlay spec>[tikz picture/node options]{<content>}
% 
% draws content boxed in a nice box
% 
\newcommand<>{\NormalBox}[2][]{%
  \only#3{\tikz[#1, every
    node/.style={shape=rectangle,draw,fill=white, drop shadow,
      #1}]\node []{#2};}
}
% **
% \OrangeBox<overlay spec>[tikz picture/node options]{<content>}
%
% draws content boxed in an orange call-out box
% 
\newcommand<>{\OrangeBox}[2][]{%
  \onslide#3{\NormalBox[fill=orange!30,draw=black!30,rounded
      corners=4pt,#1]{#2}}%
} 

%**
% \OrangeBox<overlay spec>[tikz picture/node options]{<content>}
%
% draws content boxed in an orange call-out box
% 
\newcommand<>{\ColorBox}[2][]{%
  \onslide#3{\NormalBox[fill=grey!30,draw=black!30, very thick,rounded
    corners=4pt,#1]{#2}}%
}

\newcommand<>{\NoColorBox}[2][]{%
  \onslide#3{\NormalBox[fill=white!30,draw=white!30,#1]{#2}}%
}


% \usepackage{tikz}
% \addtobeamertemplate{headline}{}
% {% 
%   \begin{flushright}
%     \begin{tikzpicture}[remember picture,overlay]
%       \node [left ]{\includegraphics[width=2cm]{gwu}};
%     \end{tikzpicture}
%   \end{flushright}
%   \vskip -0.7cm
% }

% \pgfdeclareimage[width=1.5cm]{logo}{gwu2}
% \logo{\pgfuseimage{logo}}

% \setbeamertemplate{itemize items}[default]
%%%\setbeamertemplate{itemize items}{\textbullet}
%%%\setbeamertemplate{enumerate items}[default]
%%%\setbeamertemplate{footline}[page number]

% % % logo of GWU
% \titlegraphic{
%   \centering
%   \vspace{-0.5cm}%
%   \includegraphics[width=2cm]{gwu2}
%   \vspace{0.5cm}%
% }

%%\setbeamertemplate{navigation symbols}{}%remove navigation symbols

%\usepackage[font=small,textfont={color=blue,bf},labelformat=empty]{caption}

\usepackage{verbatim}
\institute{}
% \subject{RMRF}

\def\Composite{\textsc{Composite}}
\def\composite{\textsc{Composite}}
\def\cos{\textsc{Composite}}
\def\c3{\textsc{C$^3$}}
\def\sg{\text{SuperGlue}}
\def\sm{\textsc{State Machine}}
\def\cmon{\textsc{C'Mon}}
\def\cml{\textsc{CAML}}
\def\smc{\textsc{SMC}}
\def\id{\textit{descriptor}}
\def\ids{\textit{descriptors}}
\def\Id{\textit{Descriptor}}
\def\Ids{\textit{Descriptors}}
\def\rs{\textit{resource}}
\def\rss{\textit{resources}}
\def\Rs{\textit{Resource}}

\def\emp{\textsc{EMP}}
\def\ml{\textsc{ML}}

\def\eg{\textit{e.g.}}
\def\ie{\textit{i.e.}}

\definecolor{pblue}{rgb}{0.2,0.2,0.9}
\newcommand{\blue}[1]{{\color{pblue} #1}}
%%\definecolor{pblue}{rgb}{0.2,0.2,0.9}
\newcommand{\red}[1]{{\color{red} #1}}

\makeatletter
\newcommand\titlegraphicii[1]{\def\inserttitlegraphicii{#1}}
\titlegraphicii{}
\setbeamertemplate{title page}
{
  \vbox{}
  {\usebeamercolor[fg]{titlegraphic}\inserttitlegraphic\hfill\inserttitlegraphicii\par}
  \begin{centering}
    \begin{beamercolorbox}[sep=8pt,center]{title}
      \usebeamerfont{title}\inserttitle\par%
      \ifx\insertsubtitle\@empty%
      \else%
      \vskip0.25em%
      {\usebeamerfont{subtitle}\usebeamercolor[fg]{subtitle}\insertsubtitle\par}%
      \fi%     
    \end{beamercolorbox}%
    \vskip0.5em\par
    \begin{beamercolorbox}[sep=8pt,center]{author}
      \usebeamerfont{author}\insertauthor
    \end{beamercolorbox}
    \begin{beamercolorbox}[sep=8pt,center]{date}
      \usebeamerfont{date}\insertdate
    \end{beamercolorbox}%\vskip0.5em
    \begin{beamercolorbox}[sep=8pt,center]{institute}
      \usebeamerfont{institute}\insertinstitute
    \end{beamercolorbox}
  \end{centering}
  %\vfill
}

\makeatother
%\author{Presented by {\bf Jiguo Song}}
\title{\textbf{SuperGlue: IDL-Based, System-Level Fault Tolerance for
    Embedded Systems}}
\institute{\footnotesize Computer Science Department \\ The George Washington University}
\date{\today}
%\titlegraphicii{\includegraphics[height=1.4cm,width=2cm]{gwu2}}

\usepackage{transparent}% http://ctan.org/pkg/transparent
\setbeamercovered{transparent}
\newcommand{\hidecontent}[2][0.25]{{% \hidecontent[<transparency>]{<stuff>}
    \setbox9=\hbox{#2}% Store <stuff> in \box9 to obtain height/width
    \transparent{#1}\ooalign{\usebox9\cr\color{white}\rule{\wd9}{\ht9}\cr}}}

\setbeamercovered{invisible}


\AtBeginSection[]
{
  \begin{frame}<beamer>
    \frametitle{Outlines}
    \tableofcontents[currentsection]
  \end{frame}
}


\begin{document}

\begin{frame} {\red{TO BE REMOVED - Motivation}}
\label{sec-2-1}
Motivation (8 pages)
\begin{itemize}
\item fault tolerance is important to RTOS and embedded system
  \begin{itemize}
  \item safety-critical system
  \end{itemize}
\item examples and consequences
  \begin{itemize}
  \item Hitomi
  \item SUA
  \end{itemize}
\item more complexity
  \begin{itemize}
  \item smaller feature size
  \item more complex system
  \item OS can be affected (65\%)
  \end{itemize}
\end{itemize}
\end{frame}


\begin{frame} {\red{TO BE REMOVED - system-level FT challenges}}
Why system-level fault tolerance is challenging?
\begin{itemize}
\item application-level fault tolerance (3 pages)
\begin{itemize}
\item example
\item checkpoint, N-version\ldots{}
\end{itemize}
\item system-level fault tolerance (4 pages)
\begin{itemize}
\item scheduler can be affected (examples)
\item other services as well, such as fs, lock, memory mapping manager
\item challenges
\end{itemize}
\end{itemize}
\end{frame}

\begin{frame} {\red{TO BE REMOVED - \c3}}
  The {\sc C}omputational {\sc C}rash {\sc C}art
  \begin{itemize}
  \item introduce \c3\ at high level (1 page)
    \begin{itemize}
    \item maybe need Composite?
    \end{itemize}
  \item how \c3 works (4 pages)
    \begin{itemize}
    \item fault detection
    \item micro-reboot and bring the faulty component back to a safe state
    \item bring the faulty component back to a consistent state using
      tracked descriptor state at the component interface
    \end{itemize}
  \item explain the intelligent code at the component interface (1 page)
    \begin{itemize}
    \item explain blue regions
    \end{itemize}
  \item start talking the issue with these manually written code (5
    pages)    
    \begin{itemize}
    \item needs larger font
    \item goal: less error-prone and ease the programmer burden
    \end{itemize}
  \end{itemize}
\end{frame}

\begin{frame} {\red{TO BE REMOVED - \sg\ overview}}
Introduce \sg\ (2 pages)
\begin{itemize}
\item goal: automatically generate C3 style fault recovery code
\item idea: from high-level model to low-level mechanism
\item example: use the lock example to show the goal
  \begin{itemize}
  \item needs to change the pic
  \end{itemize}
\end{itemize}
\end{frame}

\begin{frame} {\red{TO BE REMOVED - \sg\ design}}
  Get some details about \sg\ (10 pages)
  \begin{itemize}
  \item based on 3 \textbf{observation}
    \begin{itemize}
    \item \emph{descriptor-resource model} 8 page
      \begin{itemize}
      \item use examples to illustrate the model
      \item model to IDL
      \end{itemize}
    \item \emph{descriptor state machine} 6 pages
      \begin{itemize}
      \item use example to illustrate the model
      \item SM to IDL
      \end{itemize}
    \item \emph{categorized recovery mechanisms} 5 pages
      \begin{itemize}
      \item recovery mechanisms in C3
        \begin{itemize}
        \item use some examples
        \end{itemize}
      \end{itemize}
    \end{itemize}
  \item link 2 models to categorized recovery mechanisms 5 pages
    \begin{itemize}
    \item discuss the code template and predicates
    \item maybe only show the result w/o details
    \end{itemize}
  \end{itemize}
\end{frame}

\begin{frame} {\red{TO BE REMOVED - \sg\ Result}}
\label{sec-2-1}
Show some results and experiments here
\begin{itemize}
\item LOC comparison (1 page)
\item fault injection result (needs to add more here) (1 page)
\item web server results (1 page)
\end{itemize}
\end{frame}


\begin{frame}[plain]
\maketitle
\end{frame}

%%%%%%%%%%%%%%%%%%%%%%%%%%%%%%% 
% Start the contents from here
%%%%%%%%%%%%%%%%%%%%%%%%%%%%%% 

%%%%%%%%%%%%%%%%%%%%%%%%%%%%%%%%%%%%%%%%%%%%%%%%%%%%%%%%%%%% 
% Use the helicopter crash example, or blackout example
%%%%%%%%%%%%%%%%%%%%%%%%%%%%%%%%%%%%%%%%%%%%%%%%%%%%%%%%%%%% 

\section{Motivation and Challenges}

\begin{frame} {\red{TO BE REMOVED - Motivation}}
\label{sec-2-1}
Motivation (8 pages)
\begin{itemize}
\item fault tolerance is important to RTOS and embedded system
  \begin{itemize}
  \item safety-critical system
  \end{itemize}
\item examples and consequences
  \begin{itemize}
  \item Hitomi
  \item SUA
  \end{itemize}
\item more complexity
  \begin{itemize}
  \item smaller feature size
  \item more complex system
  \item OS can be affected (65\%)
  \end{itemize}
\end{itemize}
\end{frame}


\begin{frame}
  \frametitle{Real-Time and Embedded Systems}
  \begin{figure}
    \centering
    \only<1>{\hspace*{-1cm}\includegraphics[width=1\linewidth]{figures/motivation/ubiquitous/allmotiv_0}}%
  \end{figure}
  
  \bigskip \vspace*{-0.2cm}
  \hspace*{4cm}{\tiny \url{http://www.pcmag.com/encyclopedia/term/42554/embedded-system}}
\end{frame}


\begin{frame}{Consequences of Embedded System Faults}
  Japanese X-ray astronomy satellite Hitomi lost in 2016
  \begin{center}
    \mbox{\only<1>{\hspace*{-0.72cm} \includegraphics[scale=0.18]{figures/motivation/hitomi}}}%
    \mbox{\only<2>{\hspace*{-0.72cm} \includegraphics[scale=0.18]{figures/motivation/hitomi}}}%
  \end{center}
  \bigskip   \bigskip   \bigskip 
  \begin{itemize}
  \item Financial losses of \$286 million USD
  \item ``{\em It's a scientific tragedy}'' - Richard Mushotzky, UMD
  \end{itemize}
  
  \bigskip \bigskip 

  % \hspace*{-0.2cm} {\footnotesize ``{\em It's a scientific
  %     tragedy. We'd hoped for ten years.}'' - Richard Mushotzky, UMD}
  
  \vspace*{0.3cm} \hspace*{1cm}{\tiny
    \url{http://www.nature.com/news/software-error-doomed-japanese-hitomi-spacecraft-1.19835}}

  \PutAt<2>{(8cm,3cm)}{
      \begin{center}
        \includegraphics[width=4cm]{figures/motivation/radiation}
      \end{center}
    }

\end{frame}

\begin{frame}{Consequences of embedded system faults}
  Toyota Sudden Unintended Acceleration (SUA) in 2004 -- 2010
  \begin{center}
    \mbox{\only<1>{\hspace*{-0.72cm} \includegraphics[scale=0.18]{figures/motivation/toyota}}}%
  \end{center}
  \begin{itemize}
  \item 89 deaths as of May 2010 and nearly 400 U.S. lawsuits
  \item Recall 10+ million vehicles and pay \$1.2 Billion USD fine
  \end{itemize}

  \bigskip \vspace*{0.5cm}\hspace*{2cm} {\tiny
    \url{http://www.cbsnews.com/news/toyota-unintended-acceleration-has-killed-89/}}

  %{\tiny \url{http://www.toyotaproblems.com/trends/unintended-acceleration/}}
\end{frame}

\begin{frame}{Sudden Unintended Acceleration}
  \only<1>{\includegraphics[width=2.8in]{figures/motivation/toyota/toyota_etcs-i.png}}%
  \only<1>{\includegraphics[width=1.5in]{figures/motivation/toyota/Toyota_ECMsm.jpg}}%
  \only<2>{\includegraphics[width=2.8in]{figures/motivation/toyota/toyota_etcs-i.png}}%
  \only<2>{\includegraphics[width=1.5in]{figures/motivation/toyota/Toyota_ECMsm.jpg}}%
  
  \bigskip
  Uncontrolled acceleration in Toyota Camrys
  \begin{itemize}
  \item Electronic Throttle Control System (ETCS)
  \item OSEK OS, 24 tasks, 280K LOC of C
  \item bit flip in scheduler data-structures 
        
    $\to$ reproducible 30-sec uncontrolled acceleration 
  \end{itemize}

  \bigskip
  ``{\em a bit-flip there, will have the effect of killing one of the
  tasks}''

  \hspace{2.6in}-- Bookout v Toyota

  \PutAt<2>{(5.5cm,3cm)}{
    \NormalBox[text width=6.5cm]{
      \begin{center}
        \includegraphics[width=1\textwidth]{figures/motivation/bitflip_new}
      \end{center}
    }
  }

\end{frame}

\begin{frame}
  \frametitle{Embedded faults:  Bad Now, Worse Tomorrow}
  \hspace*{0.5cm}
  \includegraphics[width=3.8in]{SuperGlueFigs/Electronics_Auto}
  \medskip
  \begin{itemize}
  \item[{\color{green}+}] more functionality
  \item[{\color{red}--}] more complexity $\to$ dependability more
    challenging
  \end{itemize}

  \hspace*{3cm} {\tiny
    \url{https://www.src.org/calendar/e003506/03-sessionii-rajkumar.pdf}}
\end{frame}

\begin{frame}
  \frametitle{Embedded faults:  Bad Now, Worse Tomorrow}
  \includegraphics[width=4.3in]{figures/latent/processtechnology/processtech}

  \bigskip
  Decreasing process sizes $\to$ 5nm
  \begin{itemize}
  \item[{\color{green}+}] faster
  \item[{\color{green}+}] less power
  \item[{\color{green}+}] smaller
  \item[{\color{red}--}] increased vulnerability to HW transient
    faults
  \end{itemize}
\end{frame}

\begin{frame} {\red{TO BE REMOVED - system-level FT challenges}}
Why system-level fault tolerance is challenging?
\begin{itemize}
\item application-level fault tolerance (3 pages)
\begin{itemize}
\item example
\item checkpoint, N-version\ldots{}
\end{itemize}
\item system-level fault tolerance (4 pages)
\begin{itemize}
\item scheduler can be affected (examples)
\item other services as well, such as fs, lock, memory mapping manager
\item challenges
\end{itemize}
\end{itemize}
\end{frame}

\begin{frame}{Application-Level Fault Tolerance}
  \begin{columns}
    \begin{column}{0.6\textwidth}
      Application fault tolerance
      \begin{itemize}
      \item example recovery techniques
        \begin{itemize}
        \item recovery blocks
        \item checkpointing
        \item re-fork
        \end{itemize}
      \item temporal redundancy
        \begin{itemize}
        \item detect fault by job completion
        \item replay execution from saved state
        \end{itemize}
      \item re-execution impacts only
        
        lower-priority tasks
      \end{itemize}

    \end{column}
    \begin{column}{0.35\textwidth}
      \vspace{-3em}
      \only<1,3>{\includegraphics[width=1\textwidth]{figures/app1}}
      \only<2>{\includegraphics[width=1\textwidth]{figures/app2}}
    \end{column}
  \end{columns}
\end{frame}


\begin{frame}{System-Level Fault Tolerance}
  \begin{columns}
    \begin{column}{0.6\textwidth}
      System-level fault tolerance
      \begin{itemize}
      \item failures in
        \begin{itemize}
        \item scheduler
        \item memory mapping manager
        \item file-systems
        \item synchronization manager
        \item ...
        \end{itemize}
      \end{itemize}
    \end{column}
    \begin{column}{0.35\textwidth}
      \begin{textblock}{40}(85,12)
        \only<1> {\includegraphics[width=1\textwidth]{figures/sys1}}
        \only<2> {\includegraphics[width=1\textwidth]{figures/sys2}}
        \only<3> {\includegraphics[width=1\textwidth]{figures/sys3}}
        \only<4> {\includegraphics[width=1\textwidth]{figures/BSOD.jpg}}
      \end{textblock}
    \end{column}
  \end{columns}

  \onslide<2->{ \bigskip System components contain state for all tasks
    \begin{itemize}
    \item failure impacts memory of {\em all} tasks
    \end{itemize}
  }
  
  \onslide<3->{ \bigskip Recovery requires resources
    \begin{itemize}
    \item processing time...to recovery scheduler
    \item memory...to recover memory mapper
    \end{itemize}
  }
\end{frame}

\section{System-Level Fault Recovery - \c3}

\begin{frame} {\red{TO BE REMOVED - \c3}}
  The {\sc C}omputational {\sc C}rash {\sc C}art
  \begin{itemize}
  \item introduce \c3\ at high level (1 page)
  \item how \c3 works (4 pages)
    \begin{itemize}
    \item fault detection
    \item micro-reboot and bring the faulty component back to a safe state
    \item bring the faulty component back to a consistent state using
      tracked descriptor state at the component interface
    \end{itemize}
  \item explain the intelligent code at the component interface (1 page)
    \begin{itemize}
    \item explain blue regions
    \end{itemize}
  \item start talking the issue with these manually written code (5
    pages)    
    \begin{itemize}
    \item needs larger font
    \item goal: less error-prone and ease the programmer burden
    \end{itemize}
  \end{itemize}
\end{frame}

\begin{frame}{\c3: The {\sc C}omputational {\sc C}rash {\sc C}art}
  \begin{columns}
    \hspace{2em}
    \begin{column}{0.55\textwidth}
      {\bf \c3: Computational Crash Cart}
      \begin{itemize}
      \item resuscitate system
      \item from system-level faults
      \item predictably
      \end{itemize}
    \end{column}
    \begin{column}{0.5\textwidth}
      \includegraphics[width=2in]{figures/defib2.jpg}
    \end{column}
  \end{columns}

  \bigskip
  Main ideas
  \begin{itemize}
  \item pervasive \blue{fault isolation} $\to$ restrict propagation
  \item efficient \blue{$\mu$-reboot} of individual components
  \item \blue{interface-driven}, application-oblivious recovery
  \item on-demand recovery $\to$ bound \blue{recovery inversion}
  %\item minimize trusted computing base for fault tolerance
  \end{itemize}

  \bigskip \bigskip \bigskip

  \setbeamercolor{block
    title}{use=structure,fg=white,bg=purple!75!black}
  \setbeamercolor{block
    body}{use=structure,fg=black,bg=grey!20!white}
  \begin{block}{\scriptsize \c3\ -- an interface-driven, predictable
      system-level fault recovery mechanism}\scriptsize{J. Song,
      J. Wittrock, and G. Parmer, ``Predictable, efficient
      system-level fault tolerance in \c3'' in RTSS, 2013}\end{block}

  % \setbeamercolor{block
  %   title}{use=structure,fg=white,bg=purple!75!black}
  % \setbeamercolor{block
  %   body}{use=structure,fg=black,bg=grey!20!white}
  % \begin{block}{\footnotesize Fault Model Assumption}\scriptsize{  
  %     \begin{itemize}
  %     \item processor single-event upset and fault is detected
  %       immediately ({\sl fail-stop})
  %     \item two consecutive faults are separated by the minimum time
  %       interval
  %     \end{itemize}
  %   }\end{block}
\end{frame}

\begin{frame}{\c3\ System-Level Fault Recovery}
  \begin{columns}
    \begin{column}{0.5\textwidth}
      Recovery sequence in \c3
      \begin{enumerate}
        \onslide<2->{
        \item fault detection
        }
        \onslide<3->{
        \item safe-state recovery
          \begin{itemize}
          \item $\mu$-reboot
          \end{itemize}
        }
        \onslide<4->{
        \item consistent-state recovery
          \begin{itemize}
          \item object state recovery
          \end{itemize}
        }
      \end{enumerate}
      \vspace{0.5in}
      \hspace{0.7in}
      \includegraphics[width=1.2in]{figures/3comp_key}
    \end{column}
    \begin{column}{0.27\textwidth}
      \begin{textblock}{28}(85,12)
        \only<1> {\includegraphics[width=1\textwidth]{figures/3comp}}
        \only<2> {\includegraphics[width=1\textwidth]{figures/3comp2}}
        \only<3> {\includegraphics[width=1\textwidth]{figures/3comp3}}
        \only<4> {\includegraphics[width=1\textwidth]{figures/3comp4}}
        \end{textblock}
    \end{column}
  \end{columns}
\end{frame}

\begin{frame}{\c3\ System-Level Fault Recovery (\red{need add blue region})}
  \begin{figure}
    \centering
    \includegraphics[width=0.9\textwidth]{SuperGlueFigs/newerc3}
  \end{figure}
\end{frame}

\begin{frame}{\c3\ Implementation -- Writing Code Manually}
  \only<1>{\includegraphics[width=0.9\linewidth]{SuperGlueFigs/motivation/code1}}%
  \only<2>{\includegraphics[width=0.9\linewidth]{SuperGlueFigs/motivation/code2}}%
  \only<3>{\includegraphics[width=0.9\linewidth]{SuperGlueFigs/motivation/code3}}%
  \only<4>{\includegraphics[width=0.9\linewidth]{SuperGlueFigs/motivation/code4}}%
  \only<5>{\includegraphics[width=0.9\linewidth]{SuperGlueFigs/motivation/code5}}%
  \only<6>{\includegraphics[width=0.9\linewidth]{SuperGlueFigs/motivation/code6}}%

  \PutAt<6>{(0.8cm,3.5cm)}{
    \ColorBox[text width=11cm]{
      \LARGE {Manually writing is ad-hoc and error-prone} \\
      \vspace{0.6cm} \LARGE {Goal $\to$ {\bf automatically generate
        \c3-style recovery code}}
    }
  }
\end{frame}

\section{\sg}

\begin{frame} {\red{TO BE REMOVED - \sg\ overview}}
Introduce \sg\ (2 pages)
\begin{itemize}
\item goal: automatically generate C3 style fault recovery code
\item idea: from high-level model to low-level mechanism
\item example: use the lock example to show the goal
  \begin{itemize}
  \item needs to change the pic
  \end{itemize}
\end{itemize}
\end{frame}


\begin{frame}{\sg\ -- IDL-based System-Level Fault Tolerance}
  \begin{columns}
    \hspace{0.6em}
    \begin{column}{0.6\textwidth}
      {\bf \sg}
      \begin{itemize}
      \item automatically generate interface-driven recovery code
      \item for system-level services
      \end{itemize}
    \end{column}
    \begin{column}{0.5\textwidth}
      \includegraphics[width=1.6in]{SuperGlueFigs/SuperGlue}
    \end{column}
  \end{columns}
  
  \bigskip \bigskip \bigskip \bigskip
  \begin{columns}
    \hspace{0em}
    \begin{column}{1.1\textwidth}
      Main ideas
      % \begin{list}{$\blacksquare$}{\leftmargin=1em \itemindent=0em}
      \begin{itemize}
      \item A system \blue{model} and \blue{interface description
          language (IDL)}
      %\item The \blue{mapping} from the model to the recovery
      %  mechanisms
      \item An \blue{IDL compiler} synthesizes recovery code from the
        model
      \end{itemize}
    \end{column}
  \end{columns}
\end{frame}

\begin{frame}{\sg\ -- IDL-based System-Level Fault Tolerance}
  \begin{textblock}{100}(14,14)
    %\only<1>{\includegraphics[width=1\linewidth]{SuperGlueFigs/design/overview_1}}%
    \only<1>{\includegraphics[width=1\linewidth]{SuperGlueFigs/design/overview_2}}%
  \end{textblock}
\end{frame}



\begin{frame} {\red{TO BE REMOVED - \sg\ design}}
  Get some details about \sg\ (10 pages)
  \begin{itemize}
  \item based on 3 \textbf{observation}
    \begin{itemize}
    \item \emph{descriptor-resource model} 8 page
      \begin{itemize}
      \item use examples to illustrate the model
      \item model to IDL
      \end{itemize}
    \item \emph{descriptor state machine} 6 pages
      \begin{itemize}
      \item use example to illustrate the model
      \item SM to IDL
      \end{itemize}
    \item \emph{categorized recovery mechanisms} 5 pages
      \begin{itemize}
      \item recovery mechanisms in C3
        \begin{itemize}
        \item use some examples
        \end{itemize}
      \end{itemize}
    \end{itemize}
  \item link 2 models to categorized recovery mechanisms 5 pages
    \begin{itemize}
    \item discuss the code template and predicates
    \item maybe only show the result w/o details
    \end{itemize}
  \end{itemize}
\end{frame}


\begin{frame}{Descriptor-Resource Model}
  \begin{textblock}{100}(14,12) \vspace{-0.5em}
    \only<1>{\includegraphics[width=1\linewidth]{SuperGlueFigs/design/design_1_1}}%
    \only<2>{\includegraphics[width=1\linewidth]{SuperGlueFigs/design/design_1_2}}%
    \only<3>{\includegraphics[width=1\linewidth]{SuperGlueFigs/design/design_1_3}}%
    \only<4>{\includegraphics[width=1\linewidth]{SuperGlueFigs/design/design_1_4}}%
    \only<5>{\includegraphics[width=1\linewidth]{SuperGlueFigs/design/design_1_5}}%
    \only<6>{\includegraphics[width=1\linewidth]{SuperGlueFigs/design/design_1_6}}%
    \only<7>{\includegraphics[width=1\linewidth]{SuperGlueFigs/design/design_1_6_1}}%
  \end{textblock}
\end{frame}
\begin{frame}{Descriptor-Resource Model $\to$ IDL}
  \begin{textblock}{100}(14,12) \vspace{-0.5em}
    \only<1>{\includegraphics[width=1\linewidth]{SuperGlueFigs/design/design_2_1}}%
  \end{textblock}
\end{frame}

\begin{frame}{Component Interface Function}
  \begin{textblock}{100}(14,12) \vspace{-0.5em}
    \only<1>{\includegraphics[width=1\linewidth]{SuperGlueFigs/design/design_1_7}}%
    \only<2>{\includegraphics[width=1\linewidth]{SuperGlueFigs/design/design_1_8}}%
    \only<3>{\includegraphics[width=1\linewidth]{SuperGlueFigs/design/design_1_9}}%
  \end{textblock}
\end{frame}
% \begin{frame}{Interface Functions $\to$ IDL}
%   \begin{textblock}{100}(14,12) \vspace{-0.5em}
%     \only<1>{\includegraphics[width=1\linewidth]{SuperGlueFigs/design/design_2_2}}%
%   \end{textblock}
% \end{frame}

\begin{frame}{Descriptor State Machine}
  \begin{textblock}{100}(14,12) \vspace{-0.5em}
    %\only<1>{\includegraphics[width=1\linewidth]{SuperGlueFigs/design/design_1_10}}%
    \only<1>{\includegraphics[width=1\linewidth]{SuperGlueFigs/design/design_1_11}}%
  \end{textblock}
\end{frame}
\begin{frame}{Descriptor State Machine $\to$ IDL}
  \begin{textblock}{100}(14,12) \vspace{-0.5em}
    \only<1>{\includegraphics[width=1\linewidth]{SuperGlueFigs/design/design_2_3}}%
  \end{textblock}
\end{frame}

\begin{frame}{IDL-based Specification}
  \begin{textblock}{100}(14,12) \vspace{-0.5em}
    \only<1>{\includegraphics[width=1\linewidth]{SuperGlueFigs/design/design_2_4}}%
  \end{textblock}
\end{frame}

\begin{frame}{Toward Generating Recovery Code}
  \begin{textblock}{100}(14,12) \vspace{-0.5em}
    \only<1>{\includegraphics[width=1\linewidth]{SuperGlueFigs/design/design_2_5}}%
    %\only<2>{\includegraphics[width=1\linewidth]{SuperGlueFigs/design/design_2_6}}%
    \only<2>{\includegraphics[width=1\linewidth]{SuperGlueFigs/design/design_2_7}}%
  \end{textblock}
\end{frame}

\begin{frame}{Interface-Driven Recovery Mechanisms}
  \only<1>{\includegraphics[width=1\linewidth]{SuperGlueFigs/design/design_3_0}}%
  \only<2>{\includegraphics[width=1\linewidth]{SuperGlueFigs/design/design_3_1}}%
  \only<3>{\includegraphics[width=1\linewidth]{SuperGlueFigs/design/design_3_2}}%
  %\only<4>{\includegraphics[width=1\linewidth]{SuperGlueFigs/design/design_3_3}}%
  \only<4>{\includegraphics[width=1\linewidth]{SuperGlueFigs/design/design_3_4}}%
  \only<5>{\includegraphics[width=1\linewidth]{SuperGlueFigs/design/design_3_4_SM}}%
\end{frame}

\begin{frame}{Model $\to$ Recovery Mechanisms}
  \only<1>{\includegraphics[width=1\linewidth]{SuperGlueFigs/design/design_3_5}}%
\end{frame}

\begin{frame}{Synthesize the Recovery Code}
  \only<1>{\includegraphics[width=1\linewidth]{SuperGlueFigs/design/design_3_6}}%
  \only<2>{\includegraphics[width=1\linewidth]{SuperGlueFigs/design/design_3_7}}%
  \only<3>{\includegraphics[width=1\linewidth]{SuperGlueFigs/design/design_3_8}}%
  \only<4>{\includegraphics[width=1\linewidth]{SuperGlueFigs/design/design_3_9}}%
\end{frame}

\begin{frame}{Generate Recovery Code}
  \begin{textblock}{100}(14,12) \vspace{-0.5em}
    \only<1>{\includegraphics[width=1\linewidth]{SuperGlueFigs/design/design_4_1}}%
    \only<2>{\includegraphics[width=1\linewidth]{SuperGlueFigs/design/design_4_2}}%
    \only<3>{\includegraphics[width=1\linewidth]{SuperGlueFigs/design/design_4_3}}%
  \end{textblock}
\end{frame}

\section{Result and Experiments}

\begin{frame} {\red{TO BE REMOVED - Result}}
\label{sec-2-1}
Show some results and experiments here
\begin{itemize}
\item LOC comparison (1 page)
\item fault injection result (needs to add more here) (1 page)
\item web server results (1 page)
\end{itemize}
\end{frame}

\begin{frame}{Code Generation Result}
  \vspace{1em}
  \only<1>{\includegraphics[width=1\linewidth]{SuperGlueFigs/loc_new}}%
\end{frame}

\begin{frame}{Fault Injection Result}
  \vspace{1em}
  \only<1>{\includegraphics[width=1\linewidth]{SuperGlueFigs/loc_new_gray}}%
  
  \PutAt<1>{(7cm,5.2cm)}{
    \ColorBox[text width=4cm]{
      \begin{center}
        \includegraphics[width=1\linewidth]{SuperGlueFigs/SWIFI}
      \end{center} 
    }
  }
\end{frame}

\begin{frame}{Web Server Evaluation with Injected Faults}
    \includegraphics[width=1\linewidth]{SuperGlueFigs/WEB}
\end{frame}


\section{Conclusion}

\begin{frame}{Conclusion  (need change to something like efficient...)}
  \sg\ -- code generation for system-level fault tolerance
  \begin{itemize}
  \item descriptor-resource \blue{model} and descriptor \blue{state
      machine}
  \item \blue{IDL-based} declarative specifications
  \item \blue{compiler} for synthesizing \c3-style recovery code
  \end{itemize}
\end{frame}

\begin{frame}
  \frametitle{Thanks}
  \begin{center}
    \vspace{1in}
    \color{black}\Huge{{? $||$ /* */}}

    \vspace{0.5in}
    \color{black}\Large{{composite.seas.gwu.edu}}

    \bigskip
    \includegraphics[width=2cm]{gwu2}
  \end{center}
\end{frame}


\ifnoextra   % this will exclude all extra sides

\else 


\begin{frame}[plain]
\end{frame}

\begin{frame}[plain]
\end{frame}






\begin{frame}
  \frametitle{What is transient fault?}
  Transient (Soft-error) Faults
  \begin{itemize}
  \item Alpha and beta particles from packaging material and/or
    neutrons from cosmic rays that: invert bit stored in SRAM cell,
    dynamic latch, or gate (register or memory)
  \item Transient faults persist for only a short duration
  \end{itemize}
\end{frame}

\begin{frame}
  \frametitle{A fault stayed in the system for 4 years!}
  \begin{figure}
    \centering
    \vspace{-2em}
    \includegraphics[width=0.8\linewidth]{figures/latent/malaysia777_1.png}
  \end{figure}
  {\tiny 08/01/2005, a Boeing Company 777-200 aircraft from Australia
    to Malaysia \\A. T. S. Bureau. In-flight upset; 240km nw perth,
    wa; boeing co 777- 200, 9m-mrg. Technical report, August 2005 }

\end{frame}

\begin{frame}{Tracking Events in \cmon\ - Synchronization}
  \begin{columns}
    \begin{column}{0.6\textwidth}
      \begin{itemize}
      \item multiple-producer/single-consumer
      \item lock is {\bf not} an option
        \begin{itemize}
        \item ``chicken or the egg''
        \end{itemize}
      \item restartable wait-free ring buffer
        \begin{itemize}
        \item roll the instruction pointer back to a 512 byte aligned
          boundary, if an event entry is not valid
        \end{itemize}
      \end{itemize}
    \end{column}
    \begin{column}{0.55\textwidth}
      \includegraphics[width=0.8\textwidth]{figures/latent/CMon2_sync1}
    \end{column}
  \end{columns}
  \setbeamercolor{block
    title}{use=structure,fg=white,bg=purple!75!black}
  \setbeamercolor{block
    body}{use=structure,fg=black,bg=grey!20!white}
  \bigskip \bigskip
  \begin{block}{Invalid event entries}\scriptsize{The number of
      wasted event entries are {\it bounded} by the maximum number of
      preemptions of threads over the monitor task’s
      period}\end{block}
\end{frame}

\begin{frame}{Tracking Events in \cmon\ - Synchronization}
  \begin{columns}
    \begin{column}{0.6\textwidth}
      \begin{itemize}
      \item multiple-producer/single-consumer
      \item lock is {\bf not} an option
        \begin{itemize}
        \item ``chicken or the egg''
        \end{itemize}
      \item restartable wait-free ring buffer
        \begin{itemize}
        \item roll the instruction pointer back to a 512 byte aligned
          boundary, if an event entry is not valid
        \end{itemize}
      \end{itemize}
    \end{column}
    \begin{column}{0.55\textwidth}
      \includegraphics[width=0.8\textwidth]{figures/latent/CMon2_sync2}
    \end{column}
  \end{columns}
  \setbeamercolor{block
    title}{use=structure,fg=white,bg=purple!75!black}
  \setbeamercolor{block
    body}{use=structure,fg=black,bg=grey!20!white}
  \bigskip \bigskip
  \begin{block}{Invalid event entries}\scriptsize{The number of
      wasted event entries are {\it bounded} by the maximum number of
      preemptions of threads over the monitor task’s
      period}\end{block}
\end{frame}


\begin{frame}{Processing Events - Data Structure}
  \begin{columns}
    \hspace{-2em}
    \begin{column}{0.65\textwidth}
      Each thread execution state
      \begin{itemize}
      \item total exec time
      \item blocking time
      \item per-component cumulative exec time
      \item per-component invocation exec time
      \item dependency due to shared resource contention
      \end{itemize}
    \end{column}
    \begin{column}{0.35\textwidth}
      \includegraphics[width=0.9\textwidth]{figures/latent/CMon5}
    \end{column}
  \end{columns}
  \bigskip \bigskip \bigskip
\end{frame}

{
  \usebackgroundtemplate{\includegraphics[width=\paperwidth]{figures/LKM}}
  \begin{frame}[plain]
  \end{frame}
}


\begin{frame}
  \frametitle{System level checkpoint overhead}
  \begin{figure}
    \centering
    \includegraphics[width=1\linewidth]{figures/freertos/checkpoint_new}
  \end{figure}
\end{frame}

\begin{frame}
  \frametitle{Overhead in \c3: no faults}
  \begin{table}[!htb]
    \large
    \begin{center}
      \begin{tabular}{|l|c|c|c|c|}
        \hline
        Component & without tracking & tracking overhead \\
        \hline\hline
        Sched & 1.45 & 0.34\\
        \hline
        MM & 0.52 & 0.09 \\
        \hline
        FS & 1.65 & 1.12\\
        \hline
      \end{tabular}
    \end{center}
%    \caption{\color{black} \footnotesize{Bookkeeping Overhead
%    ($\mu$-second)}}
    \caption{\color{black} \footnotesize{infrastructure overhead ($\mu$-seconds)}}
    \vspace{-1.4em}
  \end{table}
  \bigskip
%  $\to$ Bookkeeping overhead is reasonable

  \begin{table}[!htb]
    \large
    \begin{center}
      \begin{tabular}{|l|c|c|c|c|}
        \hline
        Component & LOC & Stub LOC \\
        \hline\hline
        Sched & 162(4.8\%) & 245\\
        \hline
        MM & 93(19.2\%) & 358\\
        \hline
        FS & 46(15.1\%) & 602\\
        \hline
      \end{tabular}
    \end{center}
%    \caption{\color{black} \footnotesize{Bookkeeping Overhead
%    ($\mu$-second)}}
    \caption{\color{black} \footnotesize{lines of code added}}
    \vspace{-1.4em}
  \end{table}

\end{frame}

\begin{frame}
  \frametitle{FASSR vs Recovery Object Number}
  \begin{figure}
    \hspace*{-0.3in} 
    \includegraphics[width=1\linewidth]{figures/RTA/RTA_SCHED_OBJ}
  \end{figure}
\end{frame}

\begin{frame}
  \frametitle{FASSR vs Fault Period}
  \begin{figure}
    \hspace*{-0.3in} 
    \includegraphics[width=1\linewidth]{figures/RTA/RTA_SCHED_FPERIOD}
  \end{figure}
\end{frame}

\begin{frame}
  \frametitle{$\mu$Benchmark in \c3}
  \begin{table}[!htb]
    \small
    \begin{center}
      \hspace*{-0.2in} 
      \begin{tabular}{|l|c|c|c|c}
        \hline
        Component & $\mu$-reboot -- mem init. & $\mu$-reboot -- exe
        init. & object recovery\\
        \hline\hline
        Sched & 7.52  & 10.15  & 0.76  \\
        \hline
        MM & 16.06  & 4.00  & 0  \\
        \hline
        FS & 6.37  & 2.66  & 5.00 \\
        \hline
      \end{tabular}
    \end{center}
    \caption{\color{black} \footnotesize{recovery cost ($\mu$-seconds)}}
  \end{table}
  \bigskip
  $\to$ Both $\mu$-reboot and object recovery are efficient
\end{frame}

\begin{frame}{System-level Fault Tolerance Challenges}
  \begin{center}
    \includegraphics[width=0.35\textwidth]{figures/sys1}
  \end{center}

  \bigskip
  Recovering a \blue{consistent} state
  \begin{itemize}
  \item assumed state of {\em system objects} is comparable in all
    parts of system
    \begin{itemize}
    \item {\em thread} at specific priority is blocked?
    \item shared {\em page} mapped at specific address? 
    \item {\em file} open at offset, opened with specific flags?
    \end{itemize}
  \end{itemize}
  $\to$ rebooting subsystem is \blue{insufficient}
\end{frame}

\begin{frame}{\cos: A Component-Based OS}
  \begin{columns}
    \begin{column}{0.75\textwidth}
      System functionality as {\em components}
      \begin{itemize}
      \item user-level, protection domains
      \item fine-grained fault isolation
      \end{itemize}
      \vspace{0.5em}
      Low-level services are components 
      \begin{itemize}
      \item kernel is small and policy-less
      \item scheduling, memory mapping, synchronization...are in
        user level
      \end{itemize}

    \end{column}
    \begin{column}{0.32\textwidth}
      \hspace{-3em}
      \includegraphics[width=1.3\textwidth, height=18em]{figures/component_graph}

      {\scriptsize Example component graph}
    \end{column}
  \end{columns}
\end{frame}

\begin{frame}{\cos: A Component-Based OS}
  \begin{columns}
    \begin{column}{0.75\textwidth}
      \textcolor{gray}{System functionality as {\em components}}
      \begin{itemize}
      \item \textcolor{gray}{user-level, protection domains}
      \item \textcolor{gray}{fine-grained fault isolation}
      \end{itemize}
      \vspace{0.5em}
      \textcolor{gray}{Low-level services are components}
      \begin{itemize}
      \item \textcolor{gray}{scheduling, memory mapping, file system, lock... are in
        user level}
      \item \textcolor{gray}{kernel is small and does not have policy}
      \end{itemize}

    \end{column}
    \begin{column}{0.32\textwidth}
      \hspace{-3em}
      \includegraphics[width=1.3\textwidth, height=18em]{figures/component_graph}

      {\scriptsize Example component graph}
    \end{column}
  \end{columns}
  \PutAt<1>{(0.3cm,3.5cm)}{
    \ColorBox[text width=12cm]{
      \begin{center}
        \huge { \begin{itemize} \item Pervasive system-level fault
            isolation\end{itemize}}
      \end{center} 
    }
  }
\end{frame}

\begin{frame}{Timing Fault Detection}
  \begin{figure}
    \centering
    \includegraphics[width=0.9\textwidth]{figures/latent/examples/example1}
  \end{figure}
  \vspace{-1.5em}
  \small
  \begin{itemize}
  \item Thread's exec time in a component is longer than the
    component's WCET $\to$ {\bf{component execution overrun}}
  \item Thread has been blocked longer than the maximum blocking time
    $\to$ {\bf{unbounded PI}}
  \item Schedule the highest-priority thread with dependency $\to$
    {\bf{schedule the wrong task}}
  \item Cycle in the dependencies $\to$ {\bf{deadlock}}
  \end{itemize}
\end{frame}


\begin{frame}{Processing Events - Activation}
  \begin{columns}
    \begin{column}{0.6\textwidth}
      \begin{itemize}
      \item Time-triggered
        \begin{itemize}
        \item periodically processes events
        %\item leaves enough slack
        \end{itemize}
      \item Buffer-triggered
        \begin{itemize}
        \item processes events when any shared buffer is full
        \end{itemize}
      \item ``first-come, first-served''
      \end{itemize}
    \end{column}
    \hspace{-1em}
    \begin{column}{0.55\textwidth}
      \begin{textblock}{60}(72,20)
        \includegraphics[width=0.9\textwidth]{figures/latent/CMon3}
      \end{textblock}
    \end{column}
  \end{columns}
\end{frame}

\begin{frame}{Design Trade-off - Monitor Period}
  \begin{figure}
    \centering
    \includegraphics[width=1\textwidth]{figures/latent/tradeoff/time_trigger}
  \end{figure}
\end{frame}

\begin{frame}{Design Trade-off - Memory Consumption}
  \begin{figure}
    \centering
    \includegraphics[width=1\textwidth]{figures/latent/tradeoff/buffer_trigger}
  \end{figure}
\end{frame}

\begin{frame}
  \frametitle{Minimum memory required for system schedulability}
  \begin{figure}
    \centering
    \hspace*{-0.3in} 
    \includegraphics[width=1.1\linewidth]{figures/latent/rta/mainmem}
  \end{figure}
\end{frame}

\begin{frame}
  \frametitle{Schedulability vs Monitoring Period}
  \begin{figure}
    \centering
    \hspace*{-0.3in} 
    \includegraphics[width=1.1\linewidth]{figures/latent/rta/mainperiod_1}
  \end{figure}
\end{frame}

\begin{frame}
  \frametitle{Schedulability vs Monitoring Period}
  \begin{figure}
    \centering
    \hspace*{-0.3in} 
    \includegraphics[width=1.1\linewidth]{figures/latent/rta/mainperiod_2}
  \end{figure}
\end{frame}

\begin{frame}
  \frametitle{FASSR vs Task Utilization}
    \hspace*{-0.3in}
    \includegraphics[width=1.1\linewidth]{figures/SCHED_UTIL100}
\end{frame}

\begin{frame}{Component Invocation Event}
  \begin{figure}
    \centering
    \includegraphics[width=0.37\textwidth]{figures/latent/cmonevt/eventlog1_new}
  \end{figure}
\end{frame}
\begin{frame}{Component Invocation Event}
  \begin{figure}
    \centering
    \includegraphics[width=0.37\textwidth]{figures/latent/cmonevt/eventlog2_new}
  \end{figure}
\end{frame}
\begin{frame}{Component Invocation Event}
  \begin{figure}
    \centering
    \includegraphics[width=0.37\textwidth]{figures/latent/cmonevt/eventlog2_2_new}
  \end{figure}
\end{frame}
\begin{frame}{Component Invocation Event}
  \begin{figure}
    \centering
    \includegraphics[width=0.37\textwidth]{figures/latent/cmonevt/eventlog2_new}
  \end{figure}
\end{frame}
\begin{frame}{Component Invocation Event}
  \begin{figure}
    \centering
    \includegraphics[width=0.37\textwidth]{figures/latent/cmonevt/eventlog3_new}
  \end{figure}
\end{frame}
\begin{frame}{Component Invocation Event}
  \begin{figure}
    \centering
    \includegraphics[width=0.37\textwidth]{figures/latent/cmonevt/eventlog4_new}
  \end{figure}
\end{frame}
\begin{frame}{Component Invocation Event}
  \begin{figure}
    \centering
    \includegraphics[width=0.37\textwidth]{figures/latent/cmonevt/eventlog5_new}
  \end{figure}
\end{frame}


\begin{frame}
  \frametitle{GM Talk}
  \begin{figure}
    \centering
    \hspace*{-0.3in} 
    \includegraphics[width=1\linewidth]{SuperGlueFigs/GMtalk}
  \end{figure}
\end{frame}


\fi   % this will include all extra sides

\end{document}
